\documentclass[12pt]{article}
\usepackage{graphicx,import}
\usepackage[svgnames]{xcolor} 
\usepackage{fancyhdr}
\usepackage{subfig}
\usepackage{hyperref}
\usepackage{enumitem}
\usepackage{cite}
\usepackage[many]{tcolorbox}
\usepackage{listings }
\usepackage[a4paper, total={6in, 8in} , bottom = 25mm , top = 25mm, headheight = 1.25cm , includehead,includefoot,heightrounded ]{geometry}
\usepackage{afterpage}
\usepackage{amssymb}
\usepackage{pdflscape}
\usepackage{gensymb}
\usepackage{textcomp}
\usepackage{tikz,pgfplots}
\usepackage{xecolor}
\usepackage{rotating}
\usepackage{pdfpages}
\usepackage[Kashida]{xepersian}
\usepackage[T1]{fontenc}
\usepackage{tikz}
\usepackage[utf8]{inputenc}
\usepackage{PTSerif} 
\usepackage{seqsplit}

\usepackage[edges]{forest}

\usepackage{listings}
\usepackage{xcolor}

\hypersetup{
	colorlinks   = true, %Colours links instead of ugly boxes
	urlcolor     = blue, %Colour for external hyperlinks
	linkcolor    = blue, %Colour of internal links
	citecolor   = red %Colour of citations
}
 
\definecolor{codegreen}{rgb}{0,0.6,0}
\definecolor{codegray}{rgb}{0.5,0.5,0.5}
\definecolor{codepurple}{rgb}{0.58,0,0.82}
\definecolor{backcolour}{rgb}{0.95,0.95,0.92}
 
\NewDocumentCommand{\codeword}{v}{
\texttt{\textcolor{blue}{#1}}
}
\lstset{language=java,keywordstyle={\bfseries \color{blue}}}

\lstdefinestyle{mystyle}{
    backgroundcolor=\color{backcolour},   
    commentstyle=\color{codegreen},
    keywordstyle=\color{magenta},
    numberstyle=\tiny\color{codegray},
    stringstyle=\color{codepurple},
    basicstyle=\ttfamily\normalsize,
    breakatwhitespace=false,         
    breaklines=true,                 
    captionpos=b,                    
    keepspaces=true,                 
    numbers=left,                    
    numbersep=5pt,                  
    showspaces=false,                
    showstringspaces=false,
    showtabs=false,                  
    tabsize=2
}

\lstset{style=mystyle}

\settextfont[Scale=1.2 ,BoldFont={Bahij Nazanin-Bold.ttf} , ItalicFont = {IRNazaninIranic.ttf}]{Bahij Nazanin-Regular.ttf}
\setlatintextfont[Scale = 1.0]{Garamond}
\DefaultMathsDigits 
\DeclareMathSizes{11}{19}{13}{9} 
%\DeclareMathSizes{12}{14.4}{8}{9}





\newenvironment{changemargin}[2]{%
\begin{list}{}{%
\setlength{\topsep}{0pt}%
\setlength{\leftmargin}{#1}%
\setlength{\rightmargin}{#2}%
\setlength{\listparindent}{\parindent}%
\setlength{\itemindent}{\parindent}%
\setlength{\parsep}{\parskip}%
}%
\item[]}{\end{list}}


\definecolor{foldercolor}{RGB}{124,166,198}

\tikzset{pics/folder/.style={code={%
    \node[inner sep=0pt, minimum size=#1](-foldericon){};
    \node[folder style, inner sep=0pt, minimum width=0.3*#1, minimum height=0.6*#1, above right, xshift=0.05*#1] at (-foldericon.west){};
    \node[folder style, inner sep=0pt, minimum size=#1] at (-foldericon.center){};}
    },
    pics/folder/.default={20pt},
    folder style/.style={draw=foldercolor!80!black,top color=foldercolor!40,bottom color=foldercolor}
}

\forestset{is file/.style={edge path'/.expanded={%
        ([xshift=\forestregister{folder indent}]!u.parent anchor) |- (.child anchor)},
        inner sep=1pt},
    this folder size/.style={edge path'/.expanded={%
        ([xshift=\forestregister{folder indent}]!u.parent anchor) |- (.child anchor) pic[solid]{folder=#1}}, inner xsep=0.6*#1},
    folder tree indent/.style={before computing xy={l=#1}},
    folder icons/.style={folder, this folder size=#1, folder tree indent=3*#1},
    folder icons/.default={12pt},
}

\begin{document}


%%% title pages
\begin{titlepage}
\begin{center}
        
\vspace*{0.7cm}

\includegraphics[width=0.4\textwidth]{sharif1.png}\\
\vspace{0.5cm}
\textbf{ \Huge{\emph ‌سیگنال‌ها و سیستم‌ها} }\\
\vspace{0.5cm}
\textbf{ \Large{ تمرین چهارم} }
\vspace{0.2cm}
       
 
      \large \textbf{دانشکده مهندسی کامپیوتر}\\\vspace{0.2cm}
    \large   دانشگاه صنعتی شریف\\\vspace{0.2cm}
       \large   ﻧﯿﻢ سال دوم 00-99 \\\vspace{0.2cm}
      \noindent\rule[1ex]{\linewidth}{1pt}
استاد:\\
    \textbf{{جناب آقای دکتر منظوری شلمانی}}


    \vspace{0.15cm}
نام و نام خانوادگی:\\

       
    \textbf{{امیرمهدی نامجو - 97107212}}
\end{center}
\end{titlepage}
%%% title pages


%%% header of pages
\newpage
\pagestyle{fancy}
\fancyhf{}
\fancyfoot{}
\cfoot{\thepage}
\chead{تمرین چهارم}
\rhead{\includegraphics[width=0.1\textwidth]{sharif.png}}
\lhead{امیرمهدی نامجو}
%%% header of pages

\KashidaOff

\section{سوال اول}


\begin{enumerate}
	\item
	$$H_1(e^{j\omega}) = \frac{1+\frac{1}{2} e^{-j\omega}}{1+\frac{1}{4} e^{-j\omega}} = \frac{1 + \frac{1}{2} \cos (\omega) - j \frac{1}{2}\sin(\omega)}{1 + \frac{1}{4} \cos(\omega) - j \frac{1}{4} \sin (\omega)}$$
	$$|H_1(e^{jw})| =  \frac{(1 + \frac{1}{2} \cos (\omega))^2 +  (\frac{1}{2}\sin(\omega))^2}{(1 + \frac{1}{4} \cos(\omega))^2 + ( \frac{1}{4} \sin (\omega))^2}$$
	$$=\frac{\frac{5}{4}+ \cos (\omega)}{\frac{17}{16} + \frac{1}{2}\cos (\omega)}$$
	
	
	$$
H_{2}\left(e^{j \omega}\right) =\frac{\frac{1}{2}+e^{-j \omega}}{1+\frac{1}{4} e^{-j \omega}} \\
=\frac{\frac{1}{2}+\cos (\omega)-j \sin (\omega)}{1+\frac{1}{4} \cos (\omega)-j \frac{1}{4} \sin (\omega)}
	$$
	
	$$
	\left|H_{2}\left(e^{j \omega}\right)\right|=\frac{\left(\frac{1}{2}+\cos (\omega)\right)^{2}+\sin ^{2}(\omega)}{\left(1+\frac{1}{4} \cos (\omega)\right)^{2}+\left(\frac{1}{4} \sin (\omega)\right)^{2}}
	$$
	
	$$=\frac{\frac{5}{4}+ \cos (\omega)}{\frac{17}{16} + \frac{1}{2}\cos (\omega)}$$
		
	$$|H_1(e^{j\omega})| = |H_2(e^{j\omega})|$$	
		
		یعنی برابرند.
		
	\item
	
	$$\angle H_1 (e^{j \omega}) = -\tan^{-1} (\frac{\frac{1}{2} \sin (\omega)}{1 + \frac{1}{2} \cos (\omega)}) - (-\tan^{-1} (\frac{\frac{1}{4} \sin (\omega)}{1 + \frac{1}{4} \cos (\omega)}))$$
	$$= -\tan^{-1} (\frac{\frac{1}{2} \sin (\omega)}{1 + \frac{1}{2} \cos (\omega)}) + \tan^{-1} (\frac{\frac{1}{4} \sin (\omega)}{1 + \frac{1}{4} \cos (\omega)}))$$
	
	$$\angle H_2 (e^{j \omega}) = -\tan^{-1} (\frac{ \sin (\omega)}{\frac{1}{2} +  \cos (\omega)}) - (-\tan^{-1} (\frac{\frac{1}{4} \sin (\omega)}{1 + \frac{1}{4} \cos (\omega)}))$$
	$$= -\tan^{-1} (\frac{ \sin (\omega)}{\frac{1}{2} +  \cos (\omega)}) +\tan^{-1} (\frac{\frac{1}{4} \sin (\omega)}{1 + \frac{1}{4} \cos (\omega)}))$$
	
	تاخیر گروه به نوعی منفی نرخ تغییرات فاز نسبت به $\omega$ است.
	
	اگر از عامل مشترک دوم در هر دو عبارت برای سادگی محاسبات صرف نظر کنیم و با $\xi(\omega)$ نمایش بدهیم، داریم:
	
	$$\tau_1 =\frac{1+ 2 \cos (\omega)}{5 + 4 \cos (\omega)} + \xi(\omega)$$
	
	$$\tau_2 = \frac{2(2 + \cos (\omega))}{5 + 4 \cos (\omega)}$$
		
		به راحتی با عدد گذاری ساده می‌توان متوجه شد که در بازه
		 $[-\pi,\pi]$
		 :
		 
		$$\tau_2 > \tau_1$$
		
		و این موضوع برای سایر تناوب‌ها هم برقرار است.
		یعنی تاخیر گروه $H_2$ بزرگتر از $H_1$ است.
		
		\item
		
$$H_1(e^{j\omega}) = \frac{1}{1+\frac{1}{4} e^{-j \omega}} + \frac{\frac{1}{2}e^{-j \omega}}{1 + \frac{1}{4} e^{-j\omega}}$$

$$h_1[n] = \mathcal{F}^{-1}(H_1(e^{j \omega})) = (-\frac{1}{4})^n u[n] + \frac{1}{2} (-\frac{1}{4})^{n-1}u[n-1]$$


$$H_2(e^{j\omega}) = \frac{1}{2} \frac{1}{1+\frac{1}{4}e^{-j\omega}}+\frac{e^{-j\omega}}{1+\frac{1}{4}e^{-j\omega}}$$

$$h_2[n] = \mathcal{F}^{-1}(H_2(e^{j \omega})) = (\frac{1}{2}(-\frac{1}{4})^n u[n] + (-\frac{1}{4})^{n-1}u[n-1]$$
\end{enumerate}


\end{document}



